\documentclass[11pt]{article}
\usepackage[a4paper,margin=26mm]{geometry}
\usepackage{amsmath,amssymb,amsthm,booktabs,hyperref}
\newtheorem{theorem}{Theorem}
\newtheorem{lemma}[theorem]{Lemma}
\newcommand{\E}{\mathbb E}
\newcommand{\Pp}{\mathbb P}
\newcommand{\CE}{C_E}
\title{A joint count--moment estimate on the left probability range}
\author{}
\date{28 October 2026}
\begin{document}
\maketitle

The third absolute moment can decrease when noise is added to a Bernoulli
law with parameter below $1/3$. Estimating that decrease separately from
the own-count loss is wasteful: both estimates involve the same two
conditional third moments. Combining them first produces the positive
reserve in Lemma~\ref{lc:joint}. This is enough to extend the noise range
from $1/500$ to $1/100$ on the whole interval $p\in[1/4,2/5]$.
No subdivision of that probability interval or Bernoulli reserve is used.

\begin{theorem}[LEFT-01]\label{lc:main}
Let $n\ge42$, $1/4\le p\le2/5$, $q=1-p$, and $Y=B+U$, where
$B$ is Bernoulli with parameter $p$ and $\E(U\mid B)=0$.
Suppose the conditional laws $U_j$ satisfy
\[
 -A\le U_0\le a_0,\qquad -a_1\le U_1\le A,\qquad
 A=7/10,\quad a_0,a_1\ge0,\quad a_0+a_1\le17/250.
\]
Put $v=pq$, $s=\E U^2$, $\lambda=ns$, $\ell=(v+s)^{-1/2}$,
$X=\ell(Y-p)$, $\rho=\E|X|^3$, and $\rho_B=(1-2v)/\sqrt v$.
Assume $0<\lambda\le1/100$ and $|H|\le46/25$.
Let $S_n$ be a sum of $n$ independent copies of $X$.
Choose $k\in\mathbb Z$ and $u$ by
\[
 np+H/\ell=k+u,\qquad -109/500\le u\le391/500,
 \qquad \delta=k-np.
\]
Such an anchor always exists, and it is interior: $1\le k\le n-1$.
For every $819/2000\le C\le41/100$, writing $F_B$ for the
distribution function of $\operatorname{Bin}(n,p)$, one has
\begin{align*}
 \Delta_C={}&\sqrt n\{\Pp(S_n\le H)-\Phi(H/\sqrt n)\}-C\rho\\
 &-\sqrt n\{F_B(k)-\Phi(\delta/\sqrt{nv})\}+C\rho_B
 <-\lambda/100.
\end{align*}
\end{theorem}

The constant $\CE=(3+\sqrt{10})/(6\sqrt{2\pi})$ lies strictly between
the two stated bounds. The external strict Bernoulli theorem therefore
excludes a nonnegative excess at target $\CE$ under these hypotheses.
At $\lambda=0$ the law is itself Bernoulli. The comparison theorem has
no contact or predecessor premise. Those premises remain necessary for
the separate acquisition of its support and threshold bounds. Reflected
events require applying the acquisition and comparison to the reflected
law; a bound on one signed branch is not an unqualified two-sided result.

\section{Condition on the actual count}
Throughout set
\[
 l=-109/500,\quad h=391/500,\quad b=1/20,\quad U_*=1/100,
 \quad v_*=3/16,\quad D=17/10.
\]
The phase interval has width one, so
$k=\lfloor np+H/\ell-l\rfloor$ is an admissible choice.
Let $W_j$ be the sum of $j$ independent copies of $U_1$ and $n-j$
independent copies of $U_0$, all independent. Write
$\pi_j=\binom njp^jq^{n-j}$, set $m=\sqrt n\pi_k$, and define
\[
 f(u)=\frac{\Phi((\delta+u)/\sqrt{n(v+s)})-
                  \Phi(\delta/\sqrt{n(v+s)})}{\pi_k}.
\]
\[
 O=\Pp(W_k>u)+f(u),\qquad
 T_{\rm in}=\sum_{i=1}^{n-k}\Pp(W_{k+i}\le u-i).
\]
Exact conditioning gives
\begin{align*}
 \Pp(S_n\le H)-F_B(k)
 ={}&-\sum_{i=0}^k\pi_{k-i}\Pp(W_{k-i}>u+i)\\
 &+\sum_{i=1}^{n-k}\pi_{k+i}\Pp(W_{k+i}\le u-i).
\end{align*}
Thus, if $\pi_{k+i}/\pi_k\le R$ on the legal incoming counts,
\begin{equation}\label{lc:master}
 \Delta_C\le-m(O-RT_{\rm in})+R_N+C(\rho_B-\rho),\qquad
 R_N=\sqrt n\{\Phi(\delta/\sqrt{nv})-
                     \Phi(\delta/\sqrt{n(v+s)})\}.
\end{equation}
All outgoing events are strict; incoming events are closed. No continuity
assumption or removal of threshold atoms enters this identity.

\section{Uniform finite normalization}
We give the constants explicitly because the small improvement in
normalization helps the left range. Here $v_*\le v\le6/25$ and
$\sqrt{v+s}<491/1000$. Consequently
\[
 |\delta|<D,\qquad |\delta|<1\ (|u|\le b),\qquad
 |\delta|<101/100\ (|u|\le1/10).
\]
The segment from $p$ to $k/n$ is contained in $(1/5,9/20)$ globally,
and in $(9/40,17/40)$ in the last two cases. This proves that $k$ is
interior. On these segments, respectively,
\[
\begin{array}{c|ccc}
 &|D'''(r\Vert p)|&|\{\log(r(1-r))\}'|&1/(12r(1-r))\\\hline
\text{global}&24&4&53/100\\
\text{near}&19&16/5&12/25
\end{array}
\]
are strict upper bounds. They follow from
$D'''(r\Vert p)=-(1-2r)/(r^2(1-r)^2)$ and the endpoint bounds for
$r(1-r)$. Finite Stirling,
\[
 \log(j!)=(j+1/2)\log j-j+\tfrac12\log(2\pi)+r_j,
 \qquad 0<r_j<1/(12j),
\]
and Taylor's theorem yield
\begin{equation}\label{lc:mass}
 m=\frac{\phi(0)}{\sqrt v}
      \exp\{-\delta^2/(2nv)+e\},\qquad n|e|<Z(D_0),
\end{equation}
where, when $|\delta|\le D_0$,
\[
 Z_g(D_0)=2D_0+53/100+4D_0^3/42,\qquad
 Z_n(D_0)=(8/5)D_0+12/25+(19/6)D_0^3/42.
\]
For example the entropy error before multiplication by $n$ is bounded
by $19D_0^3/(6n^2)$ on the near segment; the prefactor and factorial
errors are bounded by $(8/5)D_0/n$ and $12/(25n)$.
The finite factorial proof is retained in
\path{manuscript/arxiv-v1/appendices/count-comparison.tex}, line 131.

At a fixed count,
\[
 \log f'(z)=-\tfrac12\log(1+s/v)+\frac{\delta^2}{2nv}
                -\frac{(\delta+z)^2}{2n(v+s)}-e.
\]
Since $1/(2v_*)+D^2/(84v_*^2)<4$, for $|z|\le b_0$ this gives
\[
 n|\log f'(z)|\le4\lambda+
             \frac{2D_0b_0+b_0^2}{2v_*}+Z(D_0).
\]
Apply $e^{-x}\ge1-x$ and $e^x\le(1-x)^{-1}$, using
$(D_0,b_0)=(D,h),(1,b),(101/100,1/10)$ and the corresponding $Z$.
In \eqref{lc:mass}, use $\sqrt v<49/100$ and
$19943/50000<\phi(0)<2/5$. The resulting bounds are
\begin{equation}\label{lc:table}
\begin{array}{c|cc}
 &|u|\le b&l\le u\le h\\\hline
 m\text{ lower bound}&18/25&57/100\\
 f'\text{ lower bound}&47/50&2/3\\
 f'\text{ upper bound}&213/200&3/2\\
 \pi_{k+i}/\pi_k\text{ upper bound}&53/50&9/5
\end{array}
\end{equation}
We also have $f'>93/100$ for $|u|\le1/10$. All bounds hold on
the whole interval of integration defining $f$.
At a binomial mode the displacement is at most one. Equation
\eqref{lc:mass} gives $\sqrt n\max_j\pi_j<1$, proving the global
ratio bound in \eqref{lc:table}. In the inner interval the first ratio
is at most $1+p/(nv)<53/50$; all subsequent ratios are at most one,
since $p-\delta<7/5<2$.

Write $s_j=\E U_j^2$. Because $|s_1-s_0|\le4s$,
$\E W_j^2=\lambda+(j-np)(s_1-s_0)$. Hence for $t=\E W_k^2$,
\begin{equation}\label{lc:variances}
 (19/21)\lambda\le t\le(23/21)\lambda\quad(|u|\le b),\qquad
 t\le(7/6)\lambda\quad\text{globally}.
\end{equation}
For the incoming counts $\E W_{k+i}^2\le iV\lambda$, where
$V=6/5$ in the inner interval and $V=4/3$ globally.

The normal repair satisfies $R_N\le(5/2)\delta_+s$, since
$0.2/v_*^{3/2}<5/2$. The inward caps fix the signs in $|Y-p|$.
Using $U_0^3\le a_0U_0^2$, $U_1^3\ge-a_1U_1^2$, and
$(1+x)^{3/2}\le1+3x/2+3x^2/8$ for $x\ge0$ gives
\[
 (v+s)^{3/2}(\rho-\rho_B)
 \ge s\{3\min(p,q)-\tfrac32(1-2v)-17/250
                           -3(1-2v)s/(8v)\}.
\]
The expression in braces is at least $-B_-$, where
\[
 B_-=\tfrac32(1-2v_*)+17/250-\tfrac34+
                         3(1-2v_*)U_*/(8\cdot42v_*).
\]
Rational squaring verifies $(41/100)B_-/v_*^{3/2}<4/3$.
Thus the separate repairs in \eqref{lc:master} are bounded by $4s$
in the inner interval and $6s$ globally. We shall improve that bound
near the negative phase origin by combining the moments first.

\section{The joint reserve}
Set $\alpha=27/50$ and $M=5\alpha^4$. The polynomial
\[
 h_\alpha(x)=\frac{-x^4+3\alpha^2x^2+2\alpha^3x}{M}
                                  \le\mathbf1_{(0,\infty)}(x)
\]
is a strict-event minorant: its numerator is
$-x(x+\alpha)^2(x-2\alpha)$, and its positive maximum is
$(9+6\sqrt3)\alpha^4/4<5\alpha^4$.
For $r\ge0$ define the familiar separate-moment lower expression
\begin{equation}\label{lc:L}
 L(r,t)=\frac{(3\alpha^2-A^2-6r^2)t-3t^2-4Art
                        -r^4+3\alpha^2r^2+2\alpha^3r}{M}.
\end{equation}

\begin{lemma}\label{lc:joint}
Under the theorem's hypotheses, if $u=-r$ and $0\le r\le b$, then
\[
 m\Pp(W_k>-r)+C(\rho-\rho_B)>mL(r,t)+\tfrac74s.
\]
\end{lemma}
\begin{proof}
Put $N_0=n-k$, $N_1=k$, $T_j=N_js_j$, and $\mu_j=\E U_j^3$.
The convex function $x^2-\kappa x$ attains its maximum at an endpoint,
so
\[
 \E U_0^4-\kappa\mu_0\le
       \max\{A^2+\kappa A,a_0^2-\kappa a_0\}s_0,
\]
and the corresponding coefficient for $U_1$ is
$\max\{A^2-\kappa A,a_1^2+\kappa a_1\}$.
Centering and independence give
\[
 \E W_k^3=N_0\mu_0+N_1\mu_1,\quad
 \E W_k^4=\sum\E U_i^4+3t^2-3\sum(\E U_i^2)^2.
\]
The exact marginal identity is
\[
 \rho=\ell^3\{v(1-2v)+3v(s_0+s_1)+p\mu_1-q\mu_0\}.
\]
In $m\E h_\alpha(W_k+r)+C(\rho-\rho_B)$, the two cubic prices
are therefore
\[
 \kappa_0=-4r-\gamma q/N_0,\qquad
 \kappa_1=-4r+\gamma p/N_1,\qquad
 \gamma=C\ell^3M/m.
\]
The outer endpoint coefficients apply if
$\kappa_0\ge-(A-a_0)$ and $\kappa_1\le A-a_1$.
Indeed $\ell^3<25/2$, $m>18/25$, $q/N_0<1/38$, $p/N_1<1/38$,
and the exact inequality
\[
 4b+\frac{(41/100)(25/2)M}{(18/25)38}<A-17/250
\]
proves both guards. The reciprocal count bounds follow from
$|\delta|<1$, $p\ge1/4$ and $q\ge3/5$.

Consequently the fourth/cubic contribution is at least
\[
 -A^2t+4Ar(T_0-T_1)+\gamma A(qs_0+ps_1)-3t^2.
\]
Here the favorable $3\sum(\E U_i^2)^2$ was discarded. Bound
$T_0-T_1\ge-t$ and $v(s_0+s_1)\ge\min(p,q)s$. Applying the same
power inequality to the Bernoulli normalization now gives
\begin{align*}
 m\Pp(W_k>-r)+C(\rho-\rho_B)&\ge mL(r,t)\\
 &+C\ell^3s\{A+3\min(p,q)-\tfrac32(1-2v)
                         -3(1-2v)s/(8v)\}.
\end{align*}
The braces are at least
\[
 B_+=A+\tfrac34-\tfrac32(1-2v_*)-
                           3(1-2v_*)U_*/(8\cdot42v_*).
\]
Finally $(819/2000)B_+>(7/4)(491/1000)^3$ and
$\sqrt{v+s}<491/1000$. This proves the claim. Notice the reserve
is proportional to $s=\lambda/n$, not to $\lambda$.
\end{proof}

\section{The own count at small phase}
The usual absolute-moment bounds
$|\E W_k^3|\le At$ and $\E W_k^4\le A^2t+3t^2$ are enough
outside the use of Lemma~\ref{lc:joint}.
After shifting the quartic by $\sigma r$, adding $cr$, dividing by
$t>0$, and setting $r=zt$, their lower expression is
\begin{align*}
 F_{\sigma,c}(z,t)={}&\frac{3\alpha^2-A^2-3t}{M}
 +\left(\frac{2\sigma}{5\alpha}+c-\frac{4At}{M}\right)z\\
 &+\frac{3\alpha^2t-6t^2}{M}z^2-\frac{t^3z^4}{M}.
\end{align*}
Let $T_*=(23/21)U_*$ and $z_0=3/2$. The following checks are
on entire rectangles, not a sampled mesh:
\[
 F_{1,-213/200}>1/4,\qquad F_{-1,47/50}>4/5
       \quad(0\le z\le z_0,\ 0\le t\le T_*).
\]
For the first polynomial, its $z$ derivative is at most
$2/(5\alpha)-213/200+6\alpha^2T_*z_0/M<0$.
It is concave in $t$, so its minimum is among $(z_0,0),(z_0,T_*)$.
For the second polynomial its $z$ derivative is bounded below by
$47/50-2/(5\alpha)-4AT_*/M-4T_*^3z_0^3/M>0$;
the quadratic derivative omitted here is nonnegative since
$3\alpha^2-6T_*>0$. Its minimum is therefore at $(0,T_*)$.
All three corner values exceed the stated bounds.

For negative phase beyond $r=z_0t$ use Cantelli:
\[
 O\ge G_t(r):=\frac{r^2}{t+r^2}-\frac{213}{200}r>t/2,
                         \qquad z_0t\le r\le b.
\]
To verify the last inequality, $G_t'$ first increases and then decreases.
It is positive at $z_0t$, so the minimum on this interval is at an
endpoint. At the switch,
$G_t(z_0t)/t=z_0^2/(1+z_0^2t)-(213/200)z_0>1/2$.
At $b$, $G_t(b)-t/2>0$. Both tests are decreasing in $t$ and are
checked at $T_*$; also $z_0T_*<b$.
For positive phase after $r=z_0t$, use
$O\ge f(r)\ge(47/50)r>(4/5)t$.

Combining the first quartic bound with Lemma~\ref{lc:joint} and
$R_N\le(5/2)s$ gives a repair of at most $(3/4)s$ when
$-z_0t\le u\le0$. For the remaining negative inner phases use the
Cantelli own bound $t/2$ and repair $4s$. For positive inner phases
the own bound is $4t/5$ with repair $4s$.

\section{The incoming counts, including the infinite tail}
We retain the quadratic-hinge comparison of Pinelis,
\emph{On the Bennett--Hoeffding inequality},
\href{https://arxiv.org/abs/0902.4058}{arXiv:0902.4058},
Proposition 2.8 and Corollary 2.11. For independent centered summands
bounded above by $A$, with total variance at most $V_0$, it gives
\[
 \Pp(Z\ge x)\le\frac{\E(P_\mu-j)_+^2}{(x/A+\mu-j)^2},
 \qquad \mu=V_0/A^2,\quad j<x/A+\mu,
\]
where $P_\mu$ is Poisson. This follows by applying its moment inequality
to the quadratic hinge at $A(j-\mu)$ and then Markov's inequality;
increasing the variance upper bound is allowed. We apply it to $-W_{k+i}$,
preserving the closed event and using the same two conditional marginals.

Set $\mu=V\lambda/A^2$. For $u\le h_0<1-A$, the first incoming
count is bounded by
\[
 Q(\mu,h_0)=\frac{\mu^2(1/2+\mu/6)}{((1-h_0)/A-1+\mu)^2}.
\]
Indeed the hinge is $j=1$ and
$\E(P_\mu-1)_+^2=\mu^2-\mu+1-e^{-\mu}\le\mu^2/2+\mu^3/6$.
For integer $i\ge2$, $(P-i)_+^2\le\binom P{i+1}$ pointwise.
For $i=2$ this is $(r-1)(r-2)\ge0$ at integer $r=P-i\ge1$;
increasing $i$ increases the binomial coefficient at fixed $r$.
The Poisson numerator is consequently at most $(i\mu)^{i+1}/(i+1)!$.
For $j=2$ or $3$ put
\[
 b_j(h_0)=1/A-1-h_0/(jA),\qquad
 J_2(\mu,h_0)=\frac{\mu^3}{3(b_2(h_0)+\mu)^2(1-3\mu)},
\]
\[
 J_3(\mu,h_0)=\frac{3\mu^4}{8(b_3(h_0)+\mu)^2(1-3\mu)}.
\]
If $h_0\ge0$, $b_j(h_0)>0$ and $3\mu<1$, these bound the sum
of incoming probabilities from count $i=j$ onwards. The hinge at count
$i$ has denominator at least $i(b_j(h_0)+\mu)$, and the successive
coefficient ratio is
$[i/(i+2)](1+1/i)^i<e<3$.
This proves the geometric sum for all counts, without truncation.
In particular $b_3(h)=59/1050>0$, whereas the second hinge cannot be
used at $h$. At large positive phase the first two counts are instead
bounded by Cantelli.

All the displayed tail bounds increase with $\lambda$. Also
$Q(c\lambda,h_0)/\lambda$ increases while
$c\lambda\le(1-h_0)/A-1$, and $J_2(c\lambda,h_0)/\lambda$ increases
when $b_2(h_0)>0$. The first assertion follows by differentiating
$\mu(1/2+\mu/6)/(a+\mu)^2$; its derivative has numerator
$a/2+\mu(a/3-1/2)>0$ for $0\le\mu\le a$.
Thus these estimates are uniform down to arbitrarily small positive noise.

\section{Closing the four phase intervals}
All remaining arithmetic consists of the following rational comparisons.
Let
\[
 \mu_I=(6/5)U_*/A^2,\qquad \mu_G=(4/3)U_*/A^2,
 \quad R_I=53/50,\quad R_G=9/5.
\]
In the inner interval define
\[
 T_-=Q(\mu_I,0)+J_2(\mu_I,0),\qquad
 T_+=Q(\mu_I,b)+J_2(\mu_I,b).
\]
For the three inner cases the expression
\[
 \frac{18}{25}\left(c\frac{19}{21}-R_I T_\pm/U_*\right)
                         -a/42-1/100
\]
is strictly positive with
\[
 (c,T_\pm,a)=(1/4,T_-,3/4),\quad(1/2,T_-,4),\quad(4/5,T_+,4).
\]
The bracket itself is positive in each case. Equations
\eqref{lc:master}--\eqref{lc:variances}, the joint estimate and the
tail monotonicity therefore give $\Delta_C<-\lambda/100$ throughout
$[-b,b]$. The first of these three margins is
$314105373/15227387000>0$.

For negative outer phase set $T_G=7U_*/6$ and
\[
 B_-^{\rm out}=
 \min_{r\in\{b,-l\}}\left\{\frac{r^2}{T_G+r^2}-\frac32r\right\}
           -R_G\{Q(\mu_G,0)+J_2(\mu_G,0)\}.
\]
The minimum really is at the endpoints: the derivative of the expression
in braces first increases and then decreases and is positive at $b$.
This handles $u\in[l,-b]$.
For the positive intervals put $t_G=4U_*/3$ and
\[
 K(u)=\frac{t_G}{t_G+(1-u)^2}+
                  \frac{2t_G}{2t_G+(2-u)^2}.
\]
On $[b,1/10]$ use the bracket
\[
 B_+^{\rm near}=(93/100)b-R_G\{K(1/10)+J_3(\mu_G,1/10)\}.
\]
On $[1/10,h]$ use
\[
 B_+^{\rm far}=\min_{u\in\{1/10,h\}}
                   \{(2/3)u-R_GK(u)\}-R_GJ_3(\mu_G,h).
\]
Here the endpoint minimum is valid by concavity, since
$3(1-h)^2>t_G$ and $3(2-h)^2>2t_G$.
For each of these three positive brackets, the exact check is
\[
                 (57/100)B-6U_*/42-U_*/100>0.
\]
The respective margins are
\[
 \frac{12211588349}{224360220000},\qquad
 \frac{5230759451}{6254703000000},\qquad
 \frac{32120174181617}{2611470546230000}.
\]
They prove the required inequality because $\lambda\le U_*$.
The four closed intervals
\[
 [l,-b],\qquad [-b,b],\qquad [b,1/10],\qquad [1/10,h]
\]
cover the entire phase domain. The inner switch uses the actual variance $t$,
and both sides include its endpoint. This completes the proof.

\section*{Verification and scope}
\path{leftcount-check.py} verifies all 88 strict scalar gates with
\texttt{Fraction}, including every branch guard, positive denominator,
phase join and final margin. Machin's formula with alternating remainders
supplies rational enclosures for $\pi$, so the stated normal and $C_E$
bounds are also checked. The checker rejects missing, duplicated, reordered
and gapped phase lists. Other negative controls reject the unguarded
endpoint price, the invalid second hinge at $h$, and the attempt to close
this $1/100$ argument after discarding its joint reserve.

The unfavorable moment sign is not hypothetical: for $p=13/50$,
$U_1=0$, and $U_0=\pm1/1000$ equiprobably, rational squaring gives
$\rho<\rho_B$. This refutes deleting that correction, not LEFT-01.
The separate actual-law script checks the new coupled inequality at exact
atoms and at tiny positive noise. Such finite examples test implementation
and conventions; they do not establish the uniform theorem.

For the acquired witnesses the left-panel heat/ray strips of v11 imply
$U_0>-.6$ and $U_1<.7$, which imply the amplitude assumption used here.
Their contact, nonnegative-excess and sharp-predecessor premises remain
ancestors, as do the raw threshold and direct geometry bounds. This result
replaces the count exclusion only for $p\in[.25,.4]$ and
$0<\lambda\le.01$ at the stated target range. It neither proves a
finite base through 41 nor excludes the remaining larger-noise region.
It does not delete an existing certificate or alter a previous theorem.
\end{document}
